% ============================================================
\chapter{Quasi-metrics and Asymmetric Topology}
\label{ch:quasi_metrics}
% ============================================================

Since Fr\'echet (1906), a distance satisfies three axioms: positivity, symmetry, and the triangle inequality. But many real-world situations are intrinsically asymmetric: travel time uphill differs from downhill, the computational cost of an approximation depends on the direction of refinement, and in theoretical computer science, Smyth's distance between programs measures a rapprochement that is not reciprocal. Dropping the symmetry axiom leads to \emph{quasi-metrics}, a framework developed notably by Hans-Peter K\"unzi from the 1980s onward, revealing a rich and surprising topology: each quasi-metric generates \emph{two} topologies, one forward and one backward.

\begin{intuition}
Quasi-metrics generalize metrics by dropping the symmetry axiom: the ``distance'' from $x$
to $y$ may differ from the distance from $y$ to $x$. This asymmetry, far from being
pathological, naturally captures notions of oriented cost, refinement, and complexity in
computer science.
\end{intuition}

% ------------------------------------------------------------
\section{Fundamental Definitions}
% ------------------------------------------------------------

\begin{definition}[Quasi-metric]
A \emph{quasi-metric} on a set $X$ is a function $d\colon X \times X \to [0,+\infty]$
satisfying:
\begin{enumerate}
  \item $d(x,x) = 0$ for all $x \in X$;
  \item $d(x,z) \leq d(x,y) + d(y,z)$ for all $x,y,z \in X$ (triangle inequality).
\end{enumerate}
If additionally $d(x,y) = 0 = d(y,x) \Rightarrow x = y$, we say $d$ is
\emph{$T_0$-separated}. The pair $(X,d)$ is a \emph{quasi-metric space}.
\end{definition}

\begin{remark}
We allow $d$ to take the value $+\infty$, which is essential for certain constructions
(infinite products, complexity domains).
\end{remark}

\begin{definition}[Conjugate and symmetrization]
Let $(X,d)$ be a quasi-metric space.
\begin{itemize}
  \item The \emph{conjugate} of $d$ is $d^{-1}(x,y) = d(y,x)$.
  \item The \emph{symmetrization} of $d$ is $d^s(x,y) = \max(d(x,y), d^{-1}(x,y))$.
  \item The \emph{sum symmetrization} is $\hat{d}(x,y) = d(x,y) + d^{-1}(x,y)$.
\end{itemize}
The symmetrization $d^s$ is a metric (possibly $[0,+\infty]$-valued) when $d$ is
$T_0$-separated.
\end{definition}

% ------------------------------------------------------------
\section{Topology Associated to a Quasi-metric}
% ------------------------------------------------------------

\begin{definition}[Open balls]
For a quasi-metric space $(X,d)$, $x \in X$, and $\varepsilon > 0$:
\[
  B_d(x,\varepsilon) = \{y \in X : d(x,y) < \varepsilon\}
\]
is the \emph{forward open ball} (or \emph{direct ball}). The \emph{forward topology}
$\tau_d$ is the topology generated by these balls.
\end{definition}

\begin{proposition}
The open balls $B_d(x,\varepsilon)$ form a base for $\tau_d$. The space $(X, \tau_d)$ is
$T_0$ if and only if $d$ is $T_0$-separated.
\end{proposition}

\begin{proof}
Let $y \in B_d(x_1,\varepsilon_1) \cap B_d(x_2,\varepsilon_2)$. Set
$\delta = \min(\varepsilon_1 - d(x_1,y), \varepsilon_2 - d(x_2,y)) > 0$. For any
$z \in B_d(y,\delta)$, $d(x_i,z) \leq d(x_i,y) + d(y,z) < d(x_i,y) + \delta \leq
\varepsilon_i$. So $B_d(y,\delta) \subseteq B_d(x_1,\varepsilon_1) \cap
B_d(x_2,\varepsilon_2)$, showing the balls form a base.

For $T_0$ separation: if $\tau_d$ is $T_0$, then for $x \neq y$, there exists
$\varepsilon > 0$ with, say, $y \notin B_d(x,\varepsilon)$, giving $d(x,y) \geq
\varepsilon > 0$. Since at least one of $d(x,y)$, $d(y,x)$ is nonzero, $d$ is
$T_0$-separated.
\end{proof}

\begin{definition}[Quasi-metric specialization order]
The specialization order of $\tau_d$ is:
\[
  x \leq_d y \quad \Longleftrightarrow \quad d(x,y) = 0.
\]
\end{definition}

% ------------------------------------------------------------
\section{Fundamental Examples}
% ------------------------------------------------------------

\begin{example}[Sorgenfrey quasi-metric]
On $\R$, define:
\[
  d_S(x,y) = \begin{cases} y - x & \text{if } y \geq x, \\ 1 & \text{if } y < x.
  \end{cases}
\]
Then $\tau_{d_S}$ is the Sorgenfrey topology (base of half-open intervals $[a,b)$).
\end{example}

\begin{example}[Baire quasi-metric]
On $\N^\N$ (infinite sequences of naturals), define:
\[
  d(f,g) = \begin{cases} 0 & \text{if } f = g, \\
  2^{-\min\{n : f(n) \neq g(n)\}} & \text{if } f \neq g \text{ and } f \sqsubseteq g, \\
  1 & \text{otherwise,} \end{cases}
\]
where $f \sqsubseteq g$ means $f$ is a prefix of $g$ (in the sense that $f(k) \leq g(k)$
for the initial indices). This quasi-metric captures the idea of approximation: $d(f,g)$ is
small when $g$ extends $f$ over a long prefix.
\end{example}

\begin{example}[Complexity quasi-metric]
Let $\Sigma^*$ be the set of words over an alphabet $\Sigma$. Define:
\[
  d(u,v) = \begin{cases} 0 & \text{if } u \text{ is a prefix of } v, \\
  2^{-|u \wedge v|} & \text{otherwise,} \end{cases}
\]
where $|u \wedge v|$ is the length of the longest common prefix. This is a $T_0$-separated
quasi-metric whose forward topology is the Scott topology on (partial) words.
\end{example}

% ------------------------------------------------------------
\section{Quasi-uniformities}
% ------------------------------------------------------------

\begin{definition}[Quasi-uniformity]
A \emph{quasi-uniformity} on $X$ is a filter $\mathcal{U}$ of subsets of $X \times X$ such
that:
\begin{enumerate}
  \item For every $U \in \mathcal{U}$, the diagonal $\Delta_X \subseteq U$.
  \item For every $U \in \mathcal{U}$, there exists $V \in \mathcal{U}$ with
    $V \circ V \subseteq U$ (where $V \circ V = \{(x,z) : \exists y,\, (x,y) \in V,\,
    (y,z) \in V\}$).
\end{enumerate}
(We do \emph{not} require $U \in \mathcal{U} \Rightarrow U^{-1} \in \mathcal{U}$.)
\end{definition}

\begin{proposition}
Every quasi-metric $d$ generates a quasi-uniformity with basic entourages
$U_\varepsilon = \{(x,y) : d(x,y) < \varepsilon\}$.
\end{proposition}

\begin{definition}[Pervin quasi-uniformity]
For any topological space $(X,\tau)$, the \emph{Pervin quasi-uniformity} is generated by
the entourages:
\[
  S_U = (U \times U) \cup (U^c \times X), \quad U \in \tau.
\]
Its associated topology is $\tau$.
\end{definition}

\begin{theorem}[Quasi-uniformizability]
Every topological space is quasi-uniformizable (e.g., via the Pervin quasi-uniformity).
\end{theorem}

% ------------------------------------------------------------
\section{Completion of Quasi-metric Spaces}
% ------------------------------------------------------------

\begin{definition}[Forward Cauchy sequence]
In a quasi-metric space $(X,d)$, a sequence $(x_n)$ is a \emph{forward-Cauchy sequence} if
for every $\varepsilon > 0$, there exists $N$ such that for all $N \leq m \leq n$,
$d(x_m, x_n) < \varepsilon$.
\end{definition}

\begin{definition}[Yoneda limit]
A sequence $(x_n)$ \emph{$d$-converges} to $\ell$ (or is \emph{Yoneda-convergent}) if for
every $\varepsilon > 0$, there exists $N$ such that for all $n \geq N$,
$d(\ell, x_n) < \varepsilon$, \emph{and} for every $y \in X$ with
$\limsup_{n\to\infty} d(y, x_n) \leq r$, we have $d(y, \ell) \leq r$.
\end{definition}

\begin{definition}[Yoneda completion]
A quasi-metric space is \emph{Yoneda-complete} if every forward-Cauchy sequence has a
Yoneda limit. The \emph{Yoneda completion} of $(X,d)$ is the smallest Yoneda-complete
quasi-metric space containing $(X,d)$ as a dense subspace.
\end{definition}

\begin{theorem}[Completion and domains]
Let $(X,d)$ be a $T_0$-separated quasi-metric space. If $d$ is finite-valued and satisfies
$d(x,y) = 0 \Rightarrow x \leq y$, then the Yoneda completion of $(X,d)$ is a dcpo
equipped with a quasi-metric compatible with the Scott topology.
\end{theorem}

% ------------------------------------------------------------
\section{Continuity and Quasi-metrics}
% ------------------------------------------------------------

\begin{definition}[Non-expansive map]
A map $f\colon (X,d_X) \to (Y,d_Y)$ is \emph{non-expansive} if
$d_Y(f(x),f(y)) \leq d_X(x,y)$ for all $x,y$.
\end{definition}

\begin{proposition}
Every non-expansive map is continuous for the forward topologies. The converse is false.
\end{proposition}

\begin{theorem}[Category of quasi-metric spaces]
The category $\mathbf{QMet}$ of $T_0$-separated quasi-metric spaces and non-expansive maps
is complete and cocomplete. The forgetful functor $\mathbf{QMet} \to \mathbf{Top}_0$
preserves limits.
\end{theorem}

% ------------------------------------------------------------
\section{Co-forward Topology and the Conjugate Pair}
% ------------------------------------------------------------

\begin{definition}[Co-forward topology]
The \emph{co-forward topology} (or conjugate topology) is $\tau_{d^{-1}}$, the topology
generated by the balls of the conjugate quasi-metric $d^{-1}$.
\end{definition}

\begin{proposition}
The topologies $\tau_d$ and $\tau_{d^{-1}}$ determine a bitopological space
$(X, \tau_d, \tau_{d^{-1}})$. The join topology $\tau_d \vee \tau_{d^{-1}}$ coincides
with the topology of the symmetrization $\tau_{d^s}$.
\end{proposition}

\begin{proof}
The join topology is generated by the open balls of $d$ and of $d^{-1}$. An open ball of
$d^s$ is $B_{d^s}(x,\varepsilon) = B_d(x,\varepsilon) \cap B_{d^{-1}}(x,\varepsilon)$,
which generates the same topology.
\end{proof}

% ------------------------------------------------------------
\section{Weighted Quasi-metrics (Preview)}
% ------------------------------------------------------------

\begin{definition}[Weighted quasi-metric]
A \emph{weighted quasi-metric} (or \emph{partial quasi-metric}) on $X$ is a pair
$(d, w)$ where $d\colon X \times X \to [0,+\infty]$ is a quasi-metric and
$w\colon X \to [0,+\infty)$ is a \emph{weight} function satisfying:
\[
  d(x,y) + w(x) = d(y,x) + w(y)
\]
for all $x,y \in X$. This generalizes the notion of partial metric.
\end{definition}

\begin{remark}
Weighted quasi-metrics arise naturally in complexity analysis, where $w(x)$ measures the
``self-distance'' or ``partiality'' of $x$. They will reappear in Chapter~9.
\end{remark}

% ------------------------------------------------------------
\section{Exercises}
% ------------------------------------------------------------

\begin{exercise}
Show that the specialization order of the forward topology of a quasi-metric $d$ is given
by $x \leq y \Leftrightarrow d(x,y) = 0$.
\end{exercise}

\begin{exercise}
Let $(X,d)$ be a quasi-metric space. Show that $d$ is a metric if and only if $\tau_d$ is
$T_1$.
\end{exercise}

\begin{exercise}
Construct a quasi-metric on $\{0,1\}$ whose forward topology is the Sierpi\'nski topology.
\end{exercise}

\begin{exercise}
Show that the Pervin quasi-uniformity of a $T_0$-space is the smallest quasi-uniformity
compatible with the topology.
\end{exercise}

\begin{exercise}[Completion of $\Q$]
Let $d(x,y) = \max(y-x, 0)$ on $\Q$. Describe the forward topology, the specialization
order, and the Yoneda completion of $(\Q, d)$.
\end{exercise}

\begin{exercise}
Show that for a finite-valued quasi-metric space $(X,d)$, the following are equivalent:
\begin{enumerate}
  \item $(X,d)$ is Yoneda-complete.
  \item Every forward Cauchy filter converges.
  \item Every forward-Cauchy sequence has a Yoneda limit.
\end{enumerate}
\end{exercise}

\begin{exercise}
Let $D$ be a dcpo with the Scott topology. Construct a quasi-metric $d$ on $D$ such that
$\tau_d$ is the Scott topology and $d(x,y) = 0 \Leftrightarrow x \sqsubseteq y$.
(Hint: use a family of Scott-continuous functions into $[0,+\infty)$.)
\end{exercise}
